\documentclass{article}
\usepackage{/globals/de}
\usepackage{/globals/laws}
\usepackage{/globals/anki}
\ankiprefix{bundestag}
\ankitopic{Bundestag}
\ankishow
\begin{document}
\section{Der Bundestag}
Der Bundestag hat als einziges gewähltes Staatsorgane eine Handvoll an aufgaben, nach \art{38-48} \GG, \gesetz{BWahlG}, der \gesetz{GO BT} (Geschäftsordnung Bundestag) und dem \gesetz{AbgG}.
\ankinote{rechtsquelle}{Rechtsquelle}{Art 38 GG, GO BT, BWahlG, AbgG}

\subsection{Aufgaben}
\begin{description}
 \item[Gesetzgebungs] Neben dem Bundesrat ist der Bundestag das zentrale Gesetzgebungsorgan.
 \item[Budgetrecht] Als Gesetzgebungsorgan kann und muss der Bundestag auch den Bundeshaushalt aufstellen.
 \item[Kontrolle der Regierung] Durch kleine und große Anfragen und durch Untersuchungsausschüsse kann der Bundestag gezielt Informationen von der Bundesregierung bekommen, durch das Konstruktive Misstrauensvotum kann der Bundestag wählen die BundeskanzlerIn abzusetzen.
 \item[Legitimierung] Der Bundestag wählt, und legitimiert somit, den Bundeskanzler, den Bundespräsidenten, die Hälfte der Bundesverfassungsgerichtsrichter.
 \item[Herstellung der Öffentlichkeit] Der Bundestag nimmt den politischen Willen der Bevölkerung wahr und vermittelt diese Debatten im öffentlichen Plenum.
\end{description}
\ankinote{aufgaben}{Aufgaben}{Gesetzgebungs (Budgetrecht), Kontrolle, Legitimierung, Öffentlichkeit}
\ankinote{aufgaben-gesetz}{Aufgaben Gesetzgebungs}{Verabschieded Gesetze}
\ankinote{aufgaben-budget}{Aufgaben Budgetrecht}{Verabschieded Bundeshaushalt}
\ankinote{aufgaben-kontrolle}{Aufgaben Kontrolle}{Große, Kleine Anfragen, Untersuchungsausschuss, Konstruktives Misstrauensvotum}
\ankinote{aufgaben-legitimierung}{Aufgaben Legitimierung}{Wählt Kanzler, Präsident, hälfte BVerfG}
\ankinote{aufgaben-öffentlichkeit}{Aufgaben Öffentlichkeit}{Nimmt Willen wahr, vermittelt in Debatten}

\subsection{Zusammensetzung}
Der Bundestag wird alle 4 Jahre gewählt (\art{39} \GG), die Anzahl der Abgeordneten folgt dem \gesetz{BWahlG}.

Es gilt das Diskontinuitätsprinzip; wird ein neuer Bundestag gebildet fällt der alte und dessen unfertige Arbeit sozusagen weg; \zb Gesetzesentwürde müssen neu eingebracht werden (\gesetz{GO BT}.
\ankinote{diskontinuität}{Diskontinuitätsprinzip}{Unfertige Arbeit verfällt beim Bilden eines neuen BT}

\subsection{Wahlen}
\art{38} \GG und das BVerfG gibt mehrere Wahlrechtsgrundsätze vor.
\begin{description}
 \item[Allgemeinheit] Jeder Bürger darf wählen.
 \item[Unmittelbarkeit] Es kann keine Zwischeninstanzen geben.
 \item[Freiheit] Jede Wahl muss ohne Zwang und Druck passieren.
 \item[Gleichheit] Jede Stimme muss gleich viel Zählen.
 \item[Geheimheit] Die Stimmabgabe muss geheim passieren.
 \item[Öffentlichkeit] Die Wahlhandlung und die Stimmauszählung müssen öffentlich einsehbar passieren.
\end{description}
\ankinote{wahl}{Wahlrechtsgrundsätze}{allgemein, unmittelbar, frei, gleich, geheim, öffentlich}
\ankinote{wahl-allgemein}{Wahl allgemein}{Jeder Bürger}
\ankinote{wahl-unmittelbar}{Wahl unmittelbar}{Keine Zwischeninstanzen}
\ankinote{wahl-frei}{Wahl frei}{Kein Druck}
\ankinote{wahl-gleich}{Wahl gleich}{Jede Stimme gleiches Gewicht}
\ankinote{wahl-geheim}{Wahl geheim}{Geheime Stimmenabgabe}
\ankinote{wahl-öffentlich}{Wahl öffentlich}{Stimmauszählung öffentlich}
Diese Grundsätze in Teilen nicht vollkommen erfüllt; durch minderjährige und manche straffällige Bürger ohne Wahlrecht, durch Propaganda, durch die Briefwahl, durch die $5\%$-Hürde, \usw.

\end{document}